% Ad Familiares 12.23 --- headnote
% ad-familiares-12-23, a little after 9 October 44 BC, Rome

\textit{Cicero to Q. Cornificius (governor of Africa Vetus), from Rome a little after 9 October 44 BC ---
Perseus dateline \textit{Scr. Romae paulo post vii id. Oct. a. 710 (44)}. Cornificius's agent Tratorius has been in Rome and has briefed Cicero on the troubles of the African province; in return Cicero sends back the news that matters most --- Octavian's first move into open politics, and Antonius's departure for Brundisium to court the four Macedonian legions.}

\textit{The letter is precious as a snapshot of Cicero's mind in the weeks before the First Philippic's open break. The young Caesar is sized up with cautious enthusiasm (``there is nothing one cannot expect him to do for the sake of praise and glory'') --- still optimism, before the discounts of 43 BC. Antonius, called by the heavy irony ``our familiar friend,'' is already isolated and dangerous. Cicero ends, characteristically, with a recommendation to philosophy as the armor for what is coming --- and with the line he repeats throughout this sequence, that to bear an evil well is not to leave it unavenged.}
